\chapter{Literature Review and Related Work}
\label{chap:relatedworks}

\section{Competitor Analysis}
\label{section:competitor-analysis}

% \begin{figure}[h]
%     \centering
%     \includegraphics[width=0.5\textwidth]{examples/asana-competitive-landscape.jpg}
%     \caption{Competitive Landscape by Asana}
% \end{figure}

\subsection*{Category 1: Character Creation Systems}

\textbf{MetaHuman Creator} is a tool developed by Epic Games and integrated 
directly into Unreal Engine that allows developers to create 
photorealistic digital human characters. \cite{metahuman} It provides 
a preset-based workflow where developers blend between pre-scanned 
facial structures and refine details using built-in sliders. Characters 
produced are fully rigged and ready for animation within Unreal Engine. 
However, MetaHuman Creator is a developer-facing tool used during 
production, not a runtime system that players interact with.

\textbf{Unreal Engine's Mutable} is a runtime character customization system 
that generates dynamic skeletal meshes, materials, and textures during 
gameplay. It is designed to be integrated into games to allow players 
to customize their characters at runtime. While Mutable provides a 
flexible and powerful foundation for character customization, it 
exposes its parameter system directly to the player through 
developer-built UI, and provides no mechanism for natural language 
input.

\begin{table}[h]
\centering
\begin{tabular}{|l|c|c|c|}
\hline
\textbf{Feature} & \textbf{MetaHuman} & \textbf{Mutable} & 
\textbf{This Project} \\
\hline
Natural Language Input & \texttimes & \texttimes & \checkmark \\
Runtime Character Customization & \texttimes & \checkmark & \checkmark \\
Unreal Engine Integration & \checkmark & \checkmark & \checkmark \\
Works with Any Mesh & \texttimes & \texttimes & \checkmark \\
Drop-in Plugin & \texttimes & \texttimes & \checkmark \\
Offline Support & \texttimes & \texttimes & \checkmark \\
\hline
\end{tabular}
\caption{Comparison: Character Creation Systems}
\end{table}

%-----------------------------------------------

\subsection*{Category 2: AI-Assisted Character Creation in Games}

\textbf{Where Winds Meet} is an action RPG that features an AI-assisted 
character creation system called Smart Customization. The system 
allows players to generate a character face by either uploading 
a photograph or recording their voice, after which the game 
produces a face preset that approximates the input. Players can 
then further refine the result using the game's standard slider 
system. This is the closest existing example to the problem this 
project addresses, as it uses AI to reduce the burden of manual 
character creation. However, the system is tightly embedded within 
a single game and is not available as a reusable tool for other 
developers. It also does not support free-form natural language 
text as an input method.

\begin{table}[h]
\centering
\begin{tabular}{|l|c|c|}
\hline
\textbf{Feature} & \textbf{Where Winds Meet} & \textbf{This Project} \\
\hline
Natural Language Text Input & \texttimes & \checkmark \\
AI-Assisted Character Creation & \checkmark & \checkmark \\
Reusable Across Games & \texttimes & \checkmark \\
Unreal Engine Integration & \texttimes & \checkmark \\
Drop-in Plugin & \texttimes & \checkmark \\
Offline Support & \texttimes & \checkmark \\
Works with Any Mesh & \texttimes & \checkmark \\
Image/Voice input & \checkmark & \texttimes \\

\hline
\end{tabular}
\caption{Comparison: AI-Assisted Character Creation in Games}
\end{table}

\section{Literature Review}
\label{section:literature-review}

Brown et al. (2020) \cite{brown2020languagemodelsfewshotlearners} introduced the few-shot prompting method for Large Language Models (LLMs), demonstrated through experiments on the GPT-3 model. Their work showed that an LLM is able to perform new tasks by conditioning on a small number of input-output examples provided directly in the prompt, without any modification to the model's parameters. Their findings showed that few-shot prompting achieved performance close to or matching state-of-the-art fine-tuned models across a range of tasks. This established few-shot prompting as a cost-effective alternative to fine-tuning, particularly in settings where large amounts of labeled training data are unavailable. In this project, few-shot prompting serves as the primary mechanism for the character creation system, guiding the LLM to interpret a player's natural language description and produce character parameter values in the correct format and value range. This approach eliminates the need for a large task-specific dataset or any model fine-tuning, making it practical for deployment in both local and cloud inference settings.
